%% Generated by scripts/md_to_aamas_tex.py -- edit paper/small_paper_draft.md
%% and regenerate. Do not edit this file by hand: the Markdown is the source of
%% truth that audit_paper_numbers.py guards.
\documentclass[sigconf,anonymous]{aamas}

\usepackage{balance}
\usepackage{booktabs}
%% tabularx gives one column type X and lets LaTeX split whatever space the
%% fixed columns leave. Two hand-rolled attempts to compute column widths from
%% character counts overflowed; this removes the guessing entirely.
\usepackage{tabularx}
\usepackage{graphicx}
\usepackage{amsmath}
%% aamas.cls (like acmart) already defines \Bbbk; amssymb redefines it and
%% aborts the build. \let\Bbbk\relax is the workaround the official template
%% uses, so keep it between the two packages.
\let\Bbbk\relax
\usepackage{amssymb}
\usepackage{stfloats}
\usepackage{url}
%% artifact paths, checkpoint hashes and CLI flags are typeset in \texttt, and
%% TeX will not break a line inside \texttt by default: a single
%% "results/shoot_eval/E1_fire_oracle_dist2_3k_s60.json" then overflows its
%% column by 30-210pt. [htt] lets those tokens hyphenate, which fixes the whole
%% class of overfull boxes at once instead of shortening each path by hand.
\usepackage[htt]{hyphenat}

\setcopyright{ifaamas}
\acmConference[AAMAS '27]{Proc.\@ of the 26th International Conference on
  Autonomous Agents and Multiagent Systems (AAMAS 2027)}{May 2027}{}{ }
\copyrightyear{2027}
\acmYear{2027}
\acmDOI{}
\acmPrice{}
\acmISBN{}
\acmSubmissionID{000}

\begin{document}

\title[Diagnosing and Surgically Correcting Conservative Decision Bias in Imitation-Learned Air Combat Policies]{Diagnosing and Surgically Correcting Conservative Decision Bias in Imitation-Learned Air Combat Policies}

\author{Anonymous}
\affiliation{\institution{Anonymous submission}\city{}\country{}}
\email{anon@example.com}

\begin{abstract}
Behaviour cloning (BC) from rule-based experts is a standard bootstrap for air-combat policies, under the implicit assumption that the expert's decisions are worth copying. On a JSBSim F-16 within-visual-range 1v1 missile-engagement benchmark we test that assumption and find it fails in a specific, measurable way: the hand-designed expert launches on only \textbf{6\%} of the steps its environment permits, and BC reproduces that restraint faithfully --- a restrictive launch preference encoded in the expert rule, not an imitation defect. We introduce the \textbf{Conditional Launch Rate (CLR)} to quantify it, and \textbf{Surgical Policy Correction (SPC)}, a localized intervention that retrains \emph{only} the launch head while every other parameter --- and therefore the entire maneuver policy --- stays bit-identical. Correcting that single decision dimension recovers the performance of an oracle that fires whenever the environment permits launch: \textbf{46.50\% $\rightarrow$ 90.75\%} kill rate on 400 paired seeds (\textbf{+44.25 pp}, exact McNemar \textbf{p = 1.04e-53}, all 177 discordant seeds favouring the correction), matching the level a frozen-trajectory intervention reaches in the hardest close-range cell (\textbf{15.0\% $\rightarrow$ 86.7\%}). Under an evading target the benefit grows rather than decays (\textbf{+48.50 pp}). The measured failure mode was localized to one discrete decision dimension, and correcting only that dimension recovered oracle-level performance while preserving the original maneuver --- the oracle being a reference defined by the environment's launch mask, not a claimed optimum.
\end{abstract}

\keywords{imitation learning, behaviour cloning, expert-induced bias, causal identification, policy intervention, air combat}

\maketitle

\section{Introduction (1 page)}

\subsection{Background}

\begin{itemize}
  \item Imitation learning (IL) bootstraps control policies in domains where online trial-and-error is expensive or unsafe: autonomous driving, manipulation, air combat.
  \item BC is the simplest IL: supervised regression/classification on expert demonstrations.
  \item Dominant assumption: the expert is a \emph{sufficient} teacher; fidelity to the expert is the objective.
\end{itemize}

\subsection{Problem: expert-induced bias, amplified by closed-loop interaction}

\begin{itemize}
  \item Rule/heuristic experts are designed, not optimized. Their decision rules encode engineering judgement that is \emph{locally} plausible but never validated in the closed loop.
  \item In interactive settings (opponent/evader present), a locally reasonable restriction --- "only launch when the shot is high quality" --- changes the entire engagement trajectory, so its cost is not locally observable.
  \item BC propagates the bias faithfully, so downstream RL or evaluation inherits it as a \emph{floor}, not a \emph{fluctuation}.
  \item Existing BC evaluation reports aggregate task success only; it does not attribute the residual to a specific decision head.
  \item \textbf{Terminology, fixed once (2026-09-18):} we say \textbf{restrictive launch preference} when describing the \emph{mechanism} --- the expert's designer would call that gate \emph{precision}, and locally it is; we reserve \textbf{expert-induced decision bias} for the \emph{measured, globally costly consequence} of inheriting it. We do not assume the expert rule is erroneous in isolation; we test whether its preference remains beneficial under closed-loop interaction, which is a different question.
\end{itemize}

\subsection{Contributions}

\begin{enumerate}
  \item \textbf{C1 --- Diagnosis: expert-induced decision bias, quantified.} We introduce the \textbf{Conditional Launch Rate (CLR)}, \texttt{P(a\_fire=1 | m\_fire=1)} --- an oracle-free metric that transfers to any commit-or-not decision (launch, hand-off, emergency stop, lane change) --- and find that the rule expert launches on only ~6\% of the steps its environment permits, with BC reproducing that preference faithfully (~7\%). We localise the cause to a \textbf{designed} launch-quality gate that is strictly stricter than the environment's legality mask (\textbackslash{}S{}3.2). This is not a case of the expert being incompetent: the expert carries a \textbf{restrictive launch preference the environment does not require}, and its closed-loop cost is the decision bias this paper measures.
  \item \textbf{C2 --- Surgical Policy Correction: a localized *intervention* on the decision dimension the diagnosis identifies.} SPC modifies \textbf{only the single binary launch head} so that it fires whenever the environment permits, while \textbf{every other parameter --- and therefore the entire maneuver policy --- is provably preserved} (bit-identical heading/speed logits and action sequences, \textbackslash{}S{}3.3). The contribution is not the implementation detail of training one head; it is \emph{isolating the behavioural dimension responsible for the failure}: the maneuver is bit-identical, so the residual difference can only be attributed to the launch decision. A5/A6 (\textbackslash{}S{}4.7) then show that the benefit comes from \textbf{which parameters are allowed to move} --- not from how many, and not from where that head started.
  \item \textbf{C3 --- Causal and locality validation: two interventions, one conclusion.} Two complementary causal arguments (\textbackslash{}S{}4.3, \textbackslash{}S{}4.4), plus the A5/A6 validation that locality is the key factor (\textbackslash{}S{}4.7). \textbf{The two interventions live at different scales, and each must be cited with its own scope attached}:
\end{enumerate}

\begin{itemize}
  \item \textbf{Intervention A (cell-specific frozen-trajectory enumeration, \texttt{dist2\_3k}, 60 seeds)}: 15.0\% $\rightarrow$ 86.7\% --- \emph{the launch decision alone can explain the performance gap};
  \item \textbf{Intervention B (primary 400-seed policy intervention, seeds 20000--20399)}: 46.50\% $\rightarrow$ 90.75\% (paired, 177:0, p$\approx$1e-53) --- \emph{the correction can be internalized into the policy}.
  \item The effect survives and grows under an evading target (+44.25 $\rightarrow$ +48.50 pp), and is replicated by an independent identification strategy (\textbackslash{}S{}4.3 E6 matched 2$\times$2: +40.0 $\rightarrow$ +50.0 pp).
\end{itemize}

\textbf{Structure of the argument} (the body follows this, one level above a conventional ablation report):

\begin{figure*}[tp]
\centering
\includegraphics[width=0.82\linewidth]{framework_fig1.pdf}
\caption{Framework. A hand-designed rule expert is imitated by behaviour cloning; the inherited launch behaviour is diagnosed with CLR; and the launch head is corrected by SPC with the maneuver held provably frozen. Annotated rates are the primary-protocol values (Table 2); the frozen-trajectory figure is the enumeration of Table 3.}
\label{fig:auto1}
\end{figure*}

\section{Related Work (1 page)}

\subsection{Imitation learning and expert demonstration}

\begin{itemize}
  \item \textbf{BC and its standard pathology.} Cloning a demonstrated policy by supervised learning goes back to ALVINN \cite{pomerleau1989alvinn} and was scaled to end-to-end driving by Bojarski et al. \cite{bojarski2016end}. The canonical account of its failure mode is DAgger \cite{ross2011dagger}, which attributes the failure to \textbf{distribution shift / compounding error} between the learner's and the expert's state visitation.
  \item \textbf{Core sentence.} \emph{Existing IL methods usually assume the demonstrations are informative and near-optimal; the learner's objective is fidelity to the expert.}
  \item \textbf{The literature that does relax optimality relaxes it as a scalar.} T-REX learns a reward from ranked suboptimal demonstrations \cite{brown2019trex}; LERP models suboptimality as reward noise \cite{huo2023lerp}; IRLEED models a demonstrator's suboptimality as \emph{reward bias plus action variance} \cite{beliaev2025irleed}; IL-from-imperfection reweights supplementary suboptimal data \cite{li2023imperfect}. In every case suboptimality is a \textbf{scalar quality} to be estimated, ranked, or down-weighted.
  \item \textbf{Our gap.} \emph{What if the expert's suboptimality is not scalar but structural - a single, nameable decision gate that the environment does not require?} A scalar quality cannot recover it: the defect is invisible in task-level success statistics (\textbackslash{}S{}4.2) and costly only in closed loop. We do not down-weight the expert; we \textbf{localize} the gate and correct exactly it.
\end{itemize}

\subsection{Policy distillation and model editing}

\begin{itemize}
  \item Knowledge distillation into a student \cite{hinton2014distilling}, policy distillation \cite{rusu2016policy}, and Actor-Mimic \cite{parisotto2016actormimic} all \textbf{transfer or compress a policy globally}, with fidelity to the teacher as the objective.
  \item Parameter-efficient adaptation (LoRA \cite{hu2022lora}) and model editing (ROME \cite{meng2022rome}) do modify a small subspace while freezing the remainder - but they edit \textbf{learned knowledge} and evaluate by task success, not by attributing a residual to one behavioural dimension.
  \item \textbf{Core sentence.} \emph{Existing approaches transfer or modify policies globally; we study localized correction of a single decision dimension, and we use the freezing itself as a causal identification device.} The distinction from distillation is precise: SPC's objective is \textbf{not} to approach the teacher but to \textbf{selectively depart} from it on one head. Freezing is not an engineering convenience here - it is the identification strategy (\textbackslash{}S{}4.4).
\end{itemize}

\subsection{Autonomous air combat decision-making}

\begin{itemize}
  \item \textbf{Closest neighbours, stated plainly.} Li et al. \cite{li2024cross} learn a within-visual-range 6-DOF combat policy whose top layer decides autopilot commands \textbf{and missile launch}, with BC and PPO cross-coordinated. Li et al. \cite{li2026imitative} learn a "pursuit-lock-launch" policy with TD3 + BC and an adaptive imitation weight in Harfang3D. Both clone a launch decision from demonstrations in a WVR missile-engagement setting.
  \item \textbf{Where we differ - two ways, both load-bearing.} (i) They treat the expert as a \textbf{bootstrap to be surpassed}; we treat it as an \textbf{object of diagnosis}, and we keep the maneuver provably frozen so the residual is attributed to the launch head rather than to aggregate learning. (ii) Their success criterion is aggregate task success; we add a \textbf{conditional} criterion (CLR) that exposes failures aggregate success hides. Notably, Li et al. \cite{li2026imitative} present selective firing as a virtue - their policy "fires only when a hit was highly probable rather than spamming launches" - the \textbf{opposite prior} to the one this paper tests, and in our regime the measurement goes the other way (\textbackslash{}S{}4.3).
\end{itemize}

\section{Method (1.5 pages)}

\subsection{Task and policy}

\begin{itemize}
  \item \textbf{Task}: 1v1 \textbf{within-visual-range} air combat, JSBSim F-16 dynamics; pursuer carries 4 short-range IR missiles (AIM-9L class); initial range 2--5 km tail-chase, target heading bias 0--60$^\circ$ (default), target speed U(180,240) m/s. Max engagement range 15 km.
  \item \textbf{Observation}: 41-dim flattened = 30-dim state features + 11-dim action-mask bits. The BC encoder consumes the 30-dim state block; the mask block is applied as a separate masked-argmax step, not fed to the network.
  \item \textbf{Action space}: \texttt{MultiDiscrete([3, 5, 1, 2])} = (speed $\Delta$, heading $\Delta$, altitude, fire).
  \item \textbf{Episode}: max 1500 steps; termination = target killed / low altitude / fled altitude / target lost / ammo exhausted / timeout.
  \item \textbf{Step duration --- state it explicitly, it is a units trap.} One environment step advances the simulation by \textbf{0.2 s} (\texttt{BaseEnv.\_\allowbreak{}acmi\_\allowbreak{}time += 0.2}), i.e. the agent runs at \textbf{5 Hz}; 1500 steps $\approx$ 300 s. This is \emph{not} the JSBSim internal frame rate. Every step$\rightarrow$second conversion in this paper uses 0.2 s/step (\texttt{\textbackslash{}S{}4.3}: 67 steps = 13.4 s, 743.5 steps = 149 s). Any figure or table that appears to use 1/60 s per step is wrong by 12$\times$.
  \item \textbf{Evaluation protocol}: \textbf{deterministic masked argmax}, JSBSim physics, \texttt{difficulty\_level} as noted per experiment.
  \item \textbf{Three policies in play}:
  \item \textbf{Rule expert} --- hand-designed, stateless: every label is a pure function of the current observation (\texttt{scripts/\allowbreak{}generate\_\allowbreak{}shoot\_\allowbreak{}rule\_\allowbreak{}expert.py}), built from standard BFM heuristics (nose alignment, DLZ depth, closure sign) plus the launch-quality gate of \textbackslash{}S{}3.2. \textbf{It has no external publication to cite --- it is this project's own controller} --- and we claim only that it is \emph{a} competent teacher (\texttt{wez\_reach\_rate = 1.00}, \texttt{hit\_rate $\approx$ 0.997}, \texttt{lost\_target\_rate = 0}), not a strong or state-of-the-art one. Every headline claim concerns \emph{what imitation inherits from a given teacher}, holding the teacher fixed, so the paper does not depend on where the teacher came from.
  \item \textbf{BC round1} --- MLP encoder (256-256-128) + 4 heads, trained on 200 expert episodes, masked cross-entropy.
  \item \textbf{SPC} --- BC with only the launch head re-trained (\textbackslash{}S{}3.3).
\end{itemize}

\subsection{C1 --- Diagnostic metric and localization}

\textbf{Conditional Launch Rate (CLR)}:

\begin{itemize}
  \item \texttt{m\_fire} is the \textbf{environment's} legality mask (index 10 of the flat mask), which requires: missiles remaining > 0; launch cooldown (30 decision steps) elapsed; target alive; \texttt{ATA < 15$^\circ$}; \texttt{MIN\_\allowbreak{}ATTACK\_\allowbreak{}DISTANCE (1500 m) $\leq$ range $\leq$ dynamic max range (3-\allowbreak{}-\allowbreak{}8 km by aspect angle)}; closure < 0 (closing).
  \item \textbf{Critical property: CLR is not a claim that 100\% is optimal.} It is a descriptive quantity. The optimal value must be established separately --- which is exactly what \textbackslash{}S{}4.3 does.
\end{itemize}

\textbf{Why the expert's CLR is low --- the designed gate.} The expert does not simply follow \texttt{m\_fire}. It requires its own stricter quality predicate:

\begin{table*}[tp]
\centering
\small
\setlength{\tabcolsep}{4pt}

\label{tab:auto2}
\begin{tabular}{lrrp{0.39\textwidth}}
\toprule
Condition & Environment mask & Expert \texttt{fire\_desired} & Effect \\
\midrule
Nose alignment & \texttt{ATA < 15$^\circ$} & \texttt{ATA < 10$^\circ$} & refuses 10--15$^\circ$ shots \\
Closure & \texttt{closure < 0} & \texttt{closure < -5 m/s} & refuses marginal closure \\
DLZ position & legality only & depth $\in$ \textbf{[0.25, 0.75]} & refuses deep-in-DLZ (closer) and far-edge shots \\
Output & --- & \texttt{fire = allowed AND desired} & cooldown/inventory handled by mask \\
\bottomrule
\end{tabular}
\end{table*}

\textbf{Measured --- the gate fully determines the expert's behaviour on legal steps.} On steps where launch is legal, the expert's \texttt{fire\_desired} is \emph{identical} to its executed \texttt{action[:,3]} (560 / 9395 on the demonstration dataset), so nothing except the gate explains what it does when launch is permitted. Under the evaluated protocol, \textbf{the observed restriction is therefore explained by the designed gate} --- not by the environment and not by BC's fitting error. The measured CLR itself is Table 2; two further measurements on a different basis, and the geometry caveat that travels with them, are in Appendix E.2.

\subsection{C2 --- Surgical Policy Correction (SPC)}

\begin{itemize}
  \item \textbf{Input}: frozen BC $\theta$\_BC; BC rollouts collected over \texttt{R} episodes.
  \item \textbf{Target}: on steps where \texttt{m\_fire = 1}, command fire. All other steps are irrelevant (the mask forces \texttt{a\_fire = 0} at inference).
  \item \textbf{Procedure}:
\end{itemize}

\begin{enumerate}
  \item Roll out $\theta$\_BC; collect \texttt{(obs\_t, m\_fire,t)}.
  \item Filter to \texttt{m\_fire,t = 1}.
  \item Train \textbf{only} \texttt{action\_heads[3]} by cross-entropy toward \texttt{a\_fire = 1}. Encoder, speed head, heading head, altitude head: \texttt{requires\_grad\_(False)}.
  \item Early stop when allowed-window fire accuracy $\geq$ 0.995.
\end{enumerate}

\begin{itemize}
  \item \textbf{Verification gates} (all must PASS; implemented in \texttt{scripts/\allowbreak{}distill\_\allowbreak{}fire\_\allowbreak{}asap.py}): \textbf{G1} heading/speed logits vs $\theta$\_BC (\texttt{max diff < 1e-9}); \textbf{G2} heading/speed action sequences 100\% identical on fixed seeds; \textbf{G3} CLR $\geq$ 99.5\% over allowed windows (\texttt{P(a\_\allowbreak{}fire=1 | m\_\allowbreak{}fire=1) $\geq$ 0.995}) on held-out rollouts; \textbf{G4} no new \texttt{lost\_target} / \texttt{bad} launch introduced.
  \item \textbf{Interpretation}: G1--G2 are the causal-identification device --- the maneuver policy is provably untouched, so any performance change is attributable to the launch decision alone.
  \item \textbf{Naming, stated once:} this pipeline is called \emph{ASAP distillation} in the code (\texttt{scripts/\allowbreak{}distill\_\allowbreak{}fire\_\allowbreak{}asap.py}) and \emph{P2A} in older docs. ASAP = "As Soon As allowable by Policy", i.e. the target behaviour. The paper-facing name is \textbf{SPC}, and the mapping is recorded in Appendix A/B.
\end{itemize}

\subsection{Why isolating the launch head is the right intervention (not arbitrary)}

\begin{itemize}
  \item The launch decision is a \textbf{discrete binary commitment} with a one-step-wide opportunity in the close-range regime (\textbackslash{}S{}4.3), so its error is not smoothable by the rest of the policy.
  \item The maneuver heads are demonstrably near-ceiling: \texttt{wez\_reach\_rate = 1.0}, \texttt{hit\_\allowbreak{}rate $\approx$ 0.997-\allowbreak{}-\allowbreak{}1.000}, \texttt{lost\_target\_rate = 0} across all evaluated suites. Re-training them would buy little and destroy the clean attribution.
  \item Contrast with end-to-end PPO fine-tuning: the project's own record (\texttt{docs/summary\_phase1.md} \textbackslash{}S{}4.5) shows that once SPC reproduces the legality policy, the current reward/environment admits no launch-timing trade-off, so PPO adds no increment and carries regression risk. $\rightarrow$ report as a \emph{documented negative result}, not as a claim about PPO in general.
\end{itemize}

\section{Experiments (2 pages)}

\subsection{Setup}

\begin{itemize}
  \item JSBSim F-16 physics; 1v1 within-visual-range missile engagement; \textbf{deterministic masked argmax}.
  \item One seed protocol throughout: \textbf{seeds 20000--20399, fresh environment per episode, geometry U(0,60)}. Hardware, wall-clock and reproduction commands: Appendix D/E.
  \item Three evaluation settings: \textbf{ID} (2--5 km tail-chase), \textbf{dist2\_3k} (hard close-range cell), \textbf{target\_evasive} (\texttt{difficulty\_level = 0.3}).
  \item \textbf{All reported numbers are verified against version-controlled artifacts}; the per-number provenance record is Appendix E.
\end{itemize}

\subsection{Diagnosis: CLR reveals a restrictive launch preference}

\textbf{The expert fires on only 6.07\% of the steps its environment permits (1131/18639), and BC inherits the preference (7.26\%).} Under an evading target both barely move (\textbf{5.40\%} / \textbf{6.95\%}). The preference is a property of the expert's designed launch gate --- strictly stricter than the environment's legality condition (\textbackslash{}S{}3.2) --- and imitation transmits it faithfully; it is not an imitation defect.

\textbf{Table 2 --- Conditional Launch Rate by policy and difficulty} (fresh env, seeds 20000--20399, U(0,60), deterministic argmax):

\begin{table*}[tp]
\centering
\small
\setlength{\tabcolsep}{4pt}

\label{tab:auto3}
\begin{tabular}{lrrr}
\toprule
Policy & CLR (d=0) & CLR (d=0.3) & Legal steps (d=0 / d=0.3) \\
\midrule
Rule expert & \textbf{6.07\%} (1131/18639) & \textbf{5.40\%} (1066/19729) & 18639 / 19729 \\
BC round1 & 7.26\% (1223/16839) & 6.95\% (1194/17169) & 16839 / 17169 \\
SPC & 100.00\% (1560/1560) & 99.11\% (1563/1577) & 1560 / 1577 \\
\bottomrule
\end{tabular}
\end{table*}

Two caveats travel with this table. \textbf{(a)} SPC's 100.00\% is true \emph{by construction} --- its training target is "fire on every legal step" and gate G3 tests exactly that; it is not a finding. \textbf{(b)} The CLR denominator is policy-dependent: firing starts a 30-step cooldown that \emph{removes} legality, so a policy that fires collapses its own opportunity set (SPC $\approx$3.9 legal steps/episode, BC $\approx$42). Cross-policy CLR is therefore an ordering of tendencies, not a like-for-like exploitation fraction --- which is why the causal claim rests on \textbackslash{}S{}4.3, where the opportunity set is held constant by construction.

One cross-check on the instrument itself: adding CLR changed nothing measurable. The d=0 run that first reported CLR reproduces the kill side of the E3 pair bit-for-bit (expert 36.75\%, BC 46.50\%, paired +9.75 pp, W/L/T 63/24/313) and reproduces BC's CLR to the digit, and the d=0.3 run does the same (25.75\%, 42.25\%, +16.50 pp, 85/19/296). Expert and BC also agree closely on every basis measured --- 6.07\% vs 7.26\%, 5.40\% vs 6.95\%, and 5.96\% vs 7.37\% on the older geometry-mixed diagnostics --- so the restriction is a stable property rather than sampling noise, and BC is in no meaningful sense more restrictive than its teacher. Appendix E.3 states the scope limit on what CLR alone can establish.

\subsection{C3(a) --- Off-policy enumeration on a frozen trajectory (the key identification)}

\textbf{Design.} Take BC's heading/speed decisions as a \emph{frozen} maneuver trajectory per seed, and enumerate a family of fire policies on top of it. Because the maneuver is fixed by construction, the maneuver policy is controlled, and differences are attributable to the launch decision.

\textbf{Table 3 --- frozen-trajectory enumeration (cell \texttt{dist2\_3k}).} 60 seeds, geometry U(0,60), deterministic masked argmax, maneuver frozen per seed from the BC round-1 checkpoint (E1, regenerated 2026-09-17):

\begin{table*}[tp]
\centering
\small
\setlength{\tabcolsep}{4pt}

\label{tab:auto4}
\begin{tabular}{p{0.61\textwidth}rrr}
\toprule
Fire policy & Kill rate & Launches/ep & Reach-4-launch \\
\midrule
\textbf{\texttt{asap} --- fire whenever the launch mask permits} & \textbf{86.7\%} & 3.85 & 86.7\% \\
\texttt{delay\_30} (first legal + 30) & 38.3\% & 3.25 & 38.3\% \\
\texttt{delay\_60} & 10.0\% & 2.78 & 10.0\% \\
\texttt{dlz\_mid} (depth 0.4--0.6) & 0.0\% & 1.27 & 0.0\% \\
\texttt{dlz\_deep} (depth $\geq$ 0.6) & 0.0\% & 1.12 & 0.0\% \\
\texttt{interval\_100} & 0.0\% & 1.95 & 0.0\% \\
\textbf{frozen BC (inherited)} & \textbf{15.0\%} & 2.37 & 15.0\% \\
\bottomrule
\end{tabular}
\end{table*}

A 2026-08-06 audit of the same family on the pre-widening geometry, and what its agreement tells us about the script versus the benchmark, are in Appendix E.10.

\textbf{Reachability audit (asap arm, measured):} median first-legal step \textbf{67} (13.4 s); \textbf{cumulative legal-window length $\approx$ 4 steps per episode} (median 4.0), i.e. roughly one legal step per DLZ transit; 52/60 killed. The "$\approx$4 legal steps per episode" figure is the load-bearing structural fact: it is what makes window \emph{usage} rather than window \emph{availability} the binding constraint. Full numbers, including the 8 timeouts whose window was still open at episode end, are in Appendix E.10.

\textbf{Reading:}

\begin{itemize}
  \item The scenario is \textbf{not infeasible} --- the environment permits a 4-launch salvo in 86.7\% of episodes, and \texttt{asap} realizes 86.7\% kills.
  \item The bottleneck is the launch policy: the inherited policy uses \textbf{2.37 of 3.85} available windows, because most windows have DLZ depth < 0.25 and are rejected by \texttt{fire\_desired}.
  \item \textbf{Every} selective alternative is worse than the mask-permissive policy, and by a wide margin (best rival \texttt{delay\_30} at 38.3\%). Within the enumerated family and on this trajectory class, the mask-permissive policy \textbf{achieves the highest kill rate}; the comparison is confined to that family and is not a claim about fire policies in general.
  \item Note what this does and does not show: it establishes that the inherited launch policy is \textbf{dominated within the enumerated family, conditional on the frozen maneuver}. It does not establish that the mask-permissive policy is optimal for arbitrary maneuvers of arbitrary policies. (\textbackslash{}S{}4.4 supplies the within-policy counterpart.)
  \item \textbf{Terminology discipline.} The mask-permissive policy is \emph{not} "always fire": it fires only on steps where the environment-level launch mask is 1. The correction removes a gate the expert added \emph{on top of} the environment's legality condition; it does not relax the environment's condition itself.
\end{itemize}

\textbf{Same enumeration under an evading target --- matched 2$\times$2 (E6, seeds 20000--20059, n=60, geometry U(0,60), deterministic masked argmax).}

\texttt{target\_evasive} changes difficulty \emph{and} nothing else if we pair it with the default cell: its config is \texttt{{difficulty\_level: 0.3}} --- all other parameters default (2--5 km). Its matched d=0 counterpart is therefore cell \texttt{id\_bias30\_60} (empty config $\equiv$ defaults). Both cells are therefore the same scenario distribution, differing \textbf{only} in the target's evasion.

\begin{table*}[tp]
\centering
\small
\setlength{\tabcolsep}{4pt}

\label{tab:auto5}
\begin{tabular}{lrrrr}
\toprule
difficulty & fire policy & Kill & Launches/ep & Legal-window median \\
\midrule
d=0 & \texttt{asap} & \textbf{93.3\%} & 3.92 & 4.0 \\
d=0 & frozen BC (inherited) & \textbf{53.3\%} & 3.18 & 42.0 \\
d=0.3 & \texttt{asap} & \textbf{96.7\%} & 3.97 & 4.0 \\
d=0.3 & frozen BC (inherited) & \textbf{46.7\%} & 3.05 & 42.0 \\
 & \textbf{\texttt{asap} - inherited gap} & \textbf{+40.0 pp $\rightarrow$ +50.0 pp} &  &  \\
\bottomrule
\end{tabular}
\end{table*}

Three things follow:

\begin{enumerate}
  \item \textbf{The direction replicates under an independent method.} Here the maneuver is frozen and only the fire policy varies; in \textbackslash{}S{}4.4/\textbackslash{}S{}4.6 the fire head is re-trained and the maneuver heads are frozen. Two different identification strategies, same answer: \textbf{evasion punishes the conservative launch policy and not the aggressive one.} \texttt{asap} rises slightly (93.3 $\rightarrow$ 96.7\%) while the inherited policy falls (53.3 $\rightarrow$ 46.7\%), so the gap widens by 10.0 pp --- the same sign as the +4.25 pp widening measured by the within-policy intervention.
  \item \textbf{The opportunity-set asymmetry is a structural constant, not a difficulty effect.} The legal-window median is 4.0 for \texttt{asap} and 42.0 for the abstaining policy --- \emph{identical at both difficulties}. Firing collapses your own window (cooldown removes legality); abstaining inflates it ~10$\times$. This is the measured basis for the \textbackslash{}S{}4.2 caveat, and it is stable across the difficulty change.
  \item \textbf{Cell and difficulty were separated, not conflated.} An earlier reading of this run compared it against the \texttt{dist2\_3k} table in \textbackslash{}S{}4.3 and would have attributed part of the difference to evasion when it was actually the 2--3 km range restriction. The matched pair above removes that confound.
\end{enumerate}

\subsection{Figure 2 --- the mechanism, made visible}

\begin{figure}[tbp]
\centering
\includegraphics[width=0.82\linewidth]{mechanism_seed20007_d00.png}
\caption{Mechanism, on one illustrative episode (seed 20007, d=0). Panel A: range and angle-off, with the environment-legal launch windows shaded and both policies' launches marked. Panel B: the decisions themselves --- launch mask, BC fire, SPC fire. The state trace is bit-identical between the two policies here (max deviation 0.0e+00 m), so the whole difference sits in Panel B. Single episode, chosen for legibility; the aggregate evidence is Tables 3 and 4.}
\label{fig:auto6}
\end{figure}

The figure is an \emph{attribution} device, not a trajectory showcase: same engagement geometry, same maneuver, different launch decision. On this seed BC uses \textbf{3 of 41} legal steps (CLR 7.32\%, against the 7.26\% aggregate) and times out, while SPC uses \textbf{4 of 4} and kills the target at step 842.

\subsection{C3(b) --- SPC intervention result (primary endpoint, E7)}

\textbf{Setting (all arms identical):} d=0, deterministic masked argmax, geometry U(0,60), \textbf{seeds 20000--20399 (n=400)}, one code path (\texttt{eval\_bc\_1v1.py}), identity recorded in each file's \texttt{run\_meta} (BC \texttt{sha256 aad05b45\textbackslash{}ldots{}}, SPC \texttt{sha256 36d79bd9\textbackslash{}ldots{}}).

\begin{table*}[tp]
\centering
\small
\setlength{\tabcolsep}{4pt}

\label{tab:auto7}
\begin{tabular}{p{0.29\textwidth}rrrrrr}
\toprule
Arm & Kill rate & Wilson 95\% & Launches/ep & CLR & lost & hit \\
\midrule
Rule expert (fresh-env path, E3) & 147/400 = \textbf{36.75\%} & [32.17, 41.58] & 2.83 & --- & 0 & 0.9973 \\
\textbf{BC round1 (frozen)} & \textbf{186/400 = 46.50\%} & [41.67, 51.40] & 3.06 & 7.26\% & 0 & 0.9984 \\
\textbf{SPC (this paper)} & \textbf{363/400 = 90.75\%} & [87.51, 93.21] & 3.90 & 100.00\% & 0 & 1.0000 \\
\bottomrule
\end{tabular}
\end{table*}

\textbf{Paired BC vs expert on the same 400 seeds:} $\Delta$kill \textbf{+9.75 pp}, win/lose/tie \textbf{63 / 24 / 313}, exact McNemar \textbf{p = 3.48e-05} ($\chi^2$cc p = 4.62e-05); 72\% of the 87 discordant seeds favour BC. $\Delta$lost = 0.00, $\Delta$launches = +0.23. This is the first \emph{significance test} attached to the BC-beats-expert claim --- the sealed Phase-1 record had only a bootstrap CI ([+3.8, +11.0] pp, 500 seeds, old geometry).

\textbf{Expert numbers in this paper are fresh-env measurements.} A second evaluation path that reuses one environment across all episodes disagrees with it on identical seeds --- 32.25\% versus 36.75\% at d=0, and a 20-seed screen shows different per-seed kill vectors rather than noise. Both paths import the same rule functions at the same commanded speed, so neither the rule nor the flight controller explains the gap; the cause is environment lifecycle, since the reused-env path builds a \emph{single} \texttt{BaseEnv} while every other script evaluated here builds a fresh one per episode. The reused-env figures (32.25\%, 25.25\% at d=0.3) are therefore excluded rather than wrong: they are measured under a different episode-construction protocol, and their artifacts are renamed \texttt{EXCLUDED\_\allowbreak{}reused\_\allowbreak{}env\_\allowbreak{}expert\_\allowbreak{}*} so no table can pick them up by accident. The consistency gate is \texttt{tests/\allowbreak{}check\_\allowbreak{}expert\_\allowbreak{}path\_\allowbreak{}consistency.sh}; Appendix E.1 lists the affected scripts and artifacts.

\textbf{Paired contingency and test} (\texttt{scripts/\allowbreak{}paired\_\allowbreak{}mcnemar.py}, exact McNemar, no scipy):

\begin{table}[tbp]
\centering
\small
\setlength{\tabcolsep}{4pt}

\label{tab:auto8}
\begin{tabular}{lr}
\toprule
 & count \\
\midrule
both kill & 186 \\
neither kill & 37 \\
\textbf{only SPC kills} & \textbf{177} \\
\textbf{only BC kills} & \textbf{0} \\
discordant total & 177 \\
\bottomrule
\end{tabular}
\end{table}

\begin{itemize}
  \item \textbf{Paired difference +44.25 pp}; exact McNemar \textbf{p = 1.04e-53} ($\chi^2$ with continuity correction p = 5.97e-40).
  \item \textbf{The direction is unanimous: all 177 discordant seeds favour the correction; there is no seed where BC succeeds and SPC fails.} This is stronger than a small p-value --- it is a \emph{dominance} statement over the evaluated scenario set, and it is the cleanest single number in the paper.
  \item \texttt{lost\_target = 0} and \texttt{launch\_quality.bad = 0} in both arms --- the corrected policy does \textbf{not} buy kills with reckless launches: premium 921 / good 639 / bad 0, i.e. \textbf{every} SPC launch still scores "good-or-better" under the expert's own quality function.
  \item Heading/speed behaviour is bit-identical to BC by construction (G1 \texttt{max diff < 1e-9}; G2 identical action sequences), so the entire gain is attributable to the launch head alone.
\end{itemize}

\textbf{SPC reaches the rule oracle's level on the same seed set (E2, seeds 20000--20399, n=400).} The oracle is "frozen BC maneuver + fire on every legal step", and SPC is "bit-identical frozen maneuver heads + a launch head whose measured CLR is 100.00\%"; both therefore fire on every legal step, and their agreement is expected rather than evidence that two independent methods concur. What it does establish is that SPC captures the \textbf{entire} oracle gap --- 46.50\% $\rightarrow$ 90.75\%, i.e. 100\% of it --- rather than some fraction. The independent replication in this paper is \textbackslash{}S{}4.3(E6) against \textbackslash{}S{}4.4/\textbackslash{}S{}4.6, where the identification strategies differ in kind (off-policy enumeration versus within-policy intervention). The per-arm identities and the aggregate-only limitation are in Appendix E.

\subsection{Robustness check 1 --- widened initial geometry (heading bias) --- \emph{a benchmark check, not a contribution}}

\textbf{A benchmark-change check, not a contribution.} Widening the heading-bias range makes the benchmark easier by $\approx$ \textbf{+4.5 pp} at identical weights (87\% $\rightarrow$ 91\%, n=100), and retraining SPC under the new geometry adds a further paired \textbf{+3.25 pp} (19 vs 6 discordant pairs, robust across 4/4 bias bins) that is \emph{not} significant unpaired.

\begin{table*}[tp]
\centering
\small
\setlength{\tabcolsep}{4pt}

\label{tab:auto9}
\begin{tabular}{p{0.28\textwidth}rp{0.33\textwidth}rr}
\toprule
Comparison (d=0, new geometry U(0,60)) & n & Kill rate & Paired diff & McNemar \\
\midrule
SPC base @ U(0,60) & 400 & 362/400 = \textbf{90.50\%} (Wilson [0.872, 0.930]) & --- & --- \\
SPC retrained under U(0,60) & 400 & 375/400 = \textbf{93.75\%} (Wilson [0.909, 0.957]) & \textbf{+3.25 pp} & \textbf{p = 0.0146} \\
same pair, unpaired Fisher & 400 & --- & --- & p = 0.1146 (n.s.) \\
\bottomrule
\end{tabular}
\end{table*}

Caveats: a single seed family (42--441), d=0 only, and Bonferroni $\alpha$=0.0125 would place p = 0.0146 just outside significance. The 87\%$\rightarrow$95\% shift sometimes quoted from earlier drafts is \emph{not} the launch-head effect; the launch-head correction is \textbackslash{}S{}4.4.

\subsection{Robustness check 2 --- target evasion (\texttt{difficulty\_level = 0.3})}

Target adds S-turn \texttt{$\pm$30$^\circ$$\cdot$d$\cdot$sin(0.3t)} plus a missile-threat break-turn and a dive to \texttt{-800$\cdot$d} m (floor 2000 m). Identical protocol to \textbackslash{}S{}4.4 --- d=0 vs d=0.3 differ \textbf{only} in \texttt{--difficulty}, same seeds 20000--20399, same code path, identity recorded in each file's \texttt{run\_meta}.

\textbf{Table 4 --- main results: kill rate and paired tests.} Fresh env per episode, seeds 20000--20399, geometry U(0,60), deterministic masked argmax; \texttt{difficulty} is the only variable between the two columns' groups. CLR values live in Table 2 (\textbackslash{}S{}4.2) and are not repeated here. Every row is generated by \texttt{scripts/\allowbreak{}collect\_\allowbreak{}matrix.py} from each artifact's own \texttt{run\_meta} --- not assembled by hand.

\begin{table*}[tp]
\centering
\small
\setlength{\tabcolsep}{4pt}

\label{tab:auto10}
\begin{tabular}{p{0.39\textwidth}rrr}
\toprule
Policy & d=0 kill & d=0.3 kill & $\Delta$ (evasion cost) \\
\midrule
Rule expert (fresh env) & \textbf{36.75\%} (147/400) & \textbf{25.75\%} (103/400) & \textbf{-11.00 pp} \\
BC round1 (inherited launch policy) & 46.50\% & \textbf{42.25\%} (169/400) & -4.25 pp \\
\textbf{SPC (corrected launch head)} & 90.75\% & \textbf{90.75\%} (363/400) & \textbf{0.00} \\
\bottomrule
\end{tabular}
\end{table*}

Read against Table 2's CLR column, the table says three things at once --- the teacher is restrictive, imitation transmits it, and evasion punishes it in proportion to how restrictive the policy is (\textbf{expert -11.00 > BC -4.25 > SPC 0.00}).

\textbf{Paired tests, both against the same 400 seeds} (primary protocol: d=0 and d=0.3, geometry U(0,60), deterministic masked argmax, fresh env per episode). The expert-vs-BC comparison and the SPC-vs-BC comparison share a seed set, so both are paired by construction rather than by assumption:

\begin{table*}[tp]
\centering
\small
\setlength{\tabcolsep}{4pt}

\label{tab:auto11}
\begin{tabular}{lp{0.28\textwidth}p{0.29\textwidth}p{0.20\textwidth}}
\toprule
Paired comparison & d=0 & d=0.3 & Test \\
\midrule
\textbf{BC - expert} & \textbf{+9.75 pp} (63/24/313, discordant 87, 72\% favouring BC) & \textbf{+16.50 pp} (85/19/296, discordant 104, 82\% favouring BC) & exact McNemar \textbf{3.48e-05} / \textbf{3.79e-11} \\
\textbf{SPC - BC} & \textbf{+44.25 pp} (discordant \textbf{177 : 0}) & \textbf{+48.50 pp} (discordant \textbf{194 : 0}) & exact McNemar \textbf{1.04e-53} / \textbf{7.97e-59} \\
\bottomrule
\end{tabular}
\end{table*}

Per-seed flips under evasion: expert \textbf{+29 gained / -73 lost}, BC \textbf{+25 / -42}. Both are loss-dominated, but the expert's net is -44 seeds against BC's -17 --- the asymmetry, not just the mean, drives the ordering above.

\textbf{Three findings:}

\begin{enumerate}
  \item \textbf{Performance.} SPC's kill rate is \emph{unchanged} under evasion (90.75\% at both difficulties) while BC falls from 46.50\% to 42.25\%, so the SPC-BC gap widens from +44.25 to \textbf{+48.50 pp} and the discordant count rises from 177 to 194. The same widening appears under the independent identification of \textbackslash{}S{}4.3 (+40.0 $\rightarrow$ +50.0 pp), so it is not an artefact of the retraining procedure: one method re-trains the launch head with the maneuver frozen, the other freezes the maneuver and swaps in rule-based launch policies.
  \item \textbf{The mechanism is preserved.} CLR barely moves on either policy --- BC \textbf{7.26\% $\rightarrow$ 6.95\%}, expert \textbf{6.07\% $\rightarrow$ 5.40\%} --- so the designed gate suppresses launches to the same degree whether or not the target manoeuvres. The restriction is a fixed property of the predicate, and the imitation inherits it as one.
  \item \textbf{Interpretation.} Evasion cost is ordered by how restrictive the launch policy is: expert \textbf{-11.00 pp}, BC \textbf{-4.25 pp}, SPC \textbf{0.00 pp}, all three legs on fresh-env evaluations over the same seeds. Evasion therefore \emph{amplifies the cost of the restrictive launch preference} rather than penalising the correction. The per-seed flips say the same thing: the expert loses 73 seeds and gains 29, BC loses 42 and gains 25 --- a suppressed launch window is unrecoverable once the target turns away. This rests on a single seed family, so it is reported as a well-supported ordering rather than a law.
\end{enumerate}

Two secondary items are in Appendix E: the direct check that \texttt{difficulty=0.3} actually takes effect (368/400 episodes change length; 382/400 keep the same kill outcome, 18 flips split 9/9, so the identical 363 at both difficulties is a genuine near-cancellation rather than a no-op), and the observation that BC's paired margin over the expert itself widens under evasion (+9.75 $\rightarrow$ +16.50 pp).

\begin{figure}[tbp]
\centering
\includegraphics[width=0.90\linewidth]{robustness_fig3.pdf}
\caption{Difficulty robustness. Kill rate by policy at d=0 and d=0.3 (primary protocol, Wilson 95\textbackslash{}\% intervals). The paired SPC--BC gap widens from +44.25 to +48.50 pp; discordant counts are 177 and 194, all favouring the correction.}
\label{fig:auto12}
\end{figure}

\subsection{Ablations: why \emph{surgical}?}

Two ablations decide whether the \emph{locality} of the correction is doing the work, or whether any correction of the same objective would do. Both run at the paper's primary protocol --- 400 seeds, d=0, the same seeds as Table 4 --- so they are directly comparable with \textbackslash{}S{}4.4 rather than being screening runs.

\textbf{Table 5 --- ablations at the primary protocol} (fresh env, seeds 20000--20399, geometry U(0,60), deterministic masked argmax):

\begin{table*}[tp]
\centering
\small
\setlength{\tabcolsep}{4pt}

\label{tab:auto13}
\begin{tabular}{p{0.26\textwidth}p{0.41\textwidth}p{0.29\textwidth}}
\toprule
Intervention & Result (400 seeds, d=0) & What it tests \\
\midrule
SPC (launch head only) & \textbf{90.75\%} (363/400), CLR 100.00\% & reference \\
\textbf{A5} --- random-init launch head, encoder and maneuver heads frozen & \textbf{90.75\%} (363/400), \textbf{0 discordant seeds vs SPC}, maneuver deviation \textbf{exactly 0.00e+00} & does the correction depend on warm-starting from BC's launch head? $\rightarrow$ \textbf{no} \\
\textbf{A6} --- full-network fine-tune, every parameter trainable & \textbf{77.50\%} (310/400), \textbf{-13.25 pp} vs SPC, exact McNemar \textbf{p = 1.33e-11}, maneuver logits move by \textbf{8.51} & is \emph{freezing} merely an engineering convenience? $\rightarrow$ \textbf{no} \\
\bottomrule
\end{tabular}
\end{table*}

\textbf{A5 rules out a warm-start effect.} Random initialisation of the editable head produced behaviourally indistinguishable results under the evaluated metric, so SPC's benefit does not depend on having started near BC's launch head. The supportable statement is exactly that: A5 and SPC are \emph{behaviourally indistinguishable under the evaluated metric}, which is not the same as being proven equivalent --- no action-logit comparison is made, and a formal claim would need a pre-registered margin (\textbackslash{}S{}E.5).

\textbf{A6 shows locality is necessary rather than convenient.} Unfreezing everything makes the correction \textbf{worse} (77.50\% against SPC's 90.75\%, p = 1.33e-11) \emph{and} moves the maneuver by 8.51 logits, so the result is both less effective and no longer attributable. Full-network adaptation therefore reduces attribution and introduces maneuver changes, indicating that unrestricted updates are not equivalent to a localized intervention.

\textbf{What the pair establishes.} The observed benefit is attributable to \textbf{the choice of editable parameters} --- which decision dimension is allowed to change --- rather than to the amount of trainable parameters or the initialisation of the launch head. Protocol, budget, identity measurements, the 60-seed screening set and the A5 training log are in Appendix E.

\section{Discussion + Conclusion (0.5 page)}

\subsection{Why this is a general concern (and why it is an \emph{interactive}-systems concern)}

\begin{itemize}
  \item Rule-based experts are \textbf{designed}, not optimized; a locally plausible quality gate can be globally costly in a closed loop, and the cost is invisible in the expert's own statistics.
  \item BC is a \emph{faithful} propagator: the bias becomes a floor under every downstream method and every evaluation.
  \item The failure mode is characteristic of interacting-agent environments: the "correct" launch time depends on the trajectory the launch itself induces, so a static quality predicate cannot be validated offline.
  \item CLR is cheap, needs no oracle, and applies to any policy with a discrete commit-or-not decision (launch, hand-off, emergency stop, lane change).
\end{itemize}

\subsection{Limitations}

\begin{itemize}
  \item \textbf{1v1 WVR only}; N-vs-N extension is future work (and the thesis's core).
  \item \textbf{Single expert, single designed gate}: we identify one instance of the failure mode, not its prevalence.
  \item \textbf{Suboptimality is established conditionally}: on the frozen BC maneuver trajectory (\textbackslash{}S{}4.3) and within the frozen-heads intervention (\textbackslash{}S{}4.4). We do not claim the mask-permissive policy is optimal in general.
  \item \textbf{Statistical scope}: the +3.25 pp geometry result is a single seed family and fails a Bonferroni $\alpha$=0.0125.
  \item \textbf{Evasion is a single interpolated setting} (\texttt{difficulty = 0.3}, one scripted evasion family), not a learned or optimised opponent. "Interactive" here means a reacting scripted target, not an adversary trained against the policy.
  \item SPC requires knowing \emph{which} head to correct; automating that (e.g. per-head CLR + ablation ranking) is open.
\end{itemize}

\subsection{Future Work}

\begin{itemize}
  \item Automate head selection from per-head diagnostics.
  \item N-vs-N: per-agent CLR and per-agent surgical correction; does the bias compound with fleet size?
  \item Launch \emph{timing} when launch cost/ammo scarcity makes launching on every mask-permitted step genuinely suboptimal --- the regime where PPO becomes meaningful (project record documents that the current reward admits no such trade-off).
  \item Integrate into the hierarchical tactical architecture planned for the thesis.
\end{itemize}

\subsection{Claim statement (exact wording to use)}

\textbf{The sentence the paper exists to support (2026-09-18):} \emph{the failure was localized to a single discrete decision dimension, and correcting only that dimension recovered oracle-level performance while preserving the original maneuver policy.} Everything below is scoping, not hedging --- this sentence is the claim, the qualifiers are what keep it honest.

We do \textbf{not} claim that the mask-permissive policy is optimal in general. We claim that \textbf{a behaviorally isolated intervention reveals the inherited launch policy to be suboptimal in the studied regime}: with the maneuver policy provably frozen (\texttt{max logit diff < 1e-9}, identical action sequences), correcting the launch head alone raises the kill rate from 46.50\% to 90.75\% on 400 paired seeds (+44.25 pp, exact McNemar p = 1.04e-53, with all 177 discordant seeds favouring the correction and none favouring the original), and off-policy enumeration on a fixed maneuver trajectory shows the inherited launch policy to be dominated by the mask-permissive policy within the enumerated fire-policy family.


%% ACM-Reference-Format is the bibliography style the aamas class ships with;
%% references are excluded from the 8-page limit, so the list may run past it.
\bibliographystyle{ACM-Reference-Format}
\bibliography{references}

\balance
\end{document}
